\input{formattingHeader}

\titleandheader{Checking the OpenLCB DCC Detection Protocol}

\begin{document}
\maketitle
\thispagestyle{firststyle}

\introductionCaveats
    {https://nbviewer.org/github/openlcb/documents/blob/master/drafts/DccDetectorWN.pdf}
    {DCC Detection Protocol Working Note}.

\section{DCC Detection Protocol Procedure}

\checkProcedure{DCC Detector Protocol Checking}

There is no PIP bit allocated to this protocol.  The checks use Event
Exchange support in PIP as a prerequisite, since the DCC Detection Protocol
is built entirely on PCER messages with specially encoded Event IDs.

A node which does not self-identify in PIP that it supports
Event Exchange will be considered to have passed these checks.
\pipsetFootnote

Each detector Event ID packs three fields into the standard 8-byte Event ID:

\begin{enumerate}
\item 6 bytes: unique ID of the detector (normally the node's own Node ID).
\item 2 bits: direction/occupancy status:
    \begin{itemize}
    \item 0b00 -- unoccupied (exit)
    \item 0b01 -- occupied, forward (entry)
    \item 0b10 -- occupied, reverse (entry)
    \item 0b11 -- occupied, direction unknown (entry)
    \end{itemize}
\item 14 bits: DCC address, encoded per S-9.2.1.1 and RCN-218:
    \begin{itemize}
    \item High byte 0x38 -- DCC short address (0--127)
    \item High byte 0x39 -- DCC consist address (0--127)
    \item 0x3800 -- track empty sentinel (short address 0)
    \item Otherwise -- DCC long address
    \end{itemize}
\end{enumerate}

The auto-detection of detector events works by comparing the upper 6 bytes
of each Producer Identified Event ID against the target node's Node ID.
Events whose upper 6 bytes match are treated as detector events.

\textbf{Note:}  These do not check live detection behaviour (PCER messages
produced when a train enters or exits a monitored section), because that
would require physical train movement which cannot be triggered remotely.
The checks instead validate the Event ID encoding through the Identify Events
and Identify Producer interactions.

\subsection{Identify Detector Events}

This section checks that the DBC produces detector-format Event IDs in
response to an Identify Events Addressed message, as required by
Working Note section 2.6:
\textit{``A detector, acting as a producer, must reply to Verify Events
and Identify Producers messages with the appropriate event(s) or range.''}

The check sends an Identify Events Addressed message to the DBC.
It then examines all Producer Identified and Producer Range Identified
replies, looking for Event IDs whose upper 6 bytes match the DBC's Node ID.

\begin{enumerate}
\item That one or more Producer Identified messages are returned.
\item That those show the DBC node as their source.
\item That at least one of the returned Event IDs has its upper 6 bytes
    equal to the DBC's Node ID (confirming it is a detector).
\end{enumerate}

\subsection{Event ID Format Validation}

This section validates the encoding of all detector Event IDs discovered
in the previous step.  It checks compliance with Working Note section 2.4.

For each detector Event ID (upper 6 bytes matching the DBC's Node ID),
the check validates:

\begin{enumerate}
\item That the 2-bit direction/occupancy field contains a valid value
    (0x00 through 0x03).
\item That the upper byte of the 14-bit address field does not use a
    reserved prefix (0x3A through 0x3F).
\item That short addresses (prefix 0x38) have an address value in the
    range 0--127.
\item That consist addresses (prefix 0x39) have an address value in the
    range 0--127.
\end{enumerate}

\subsection{Identify Producer Replies}

This section checks that the DBC replies to Identify Producer messages
for each of its detector Event IDs, as required by Working Note section 2.6:
\textit{``A detector, acting as a producer, must reply to Verify Events
and Identify Producers messages with the appropriate event(s) or range.''}

The check first gathers all individual (non-range) detector Event IDs
via Identify Events Addressed.  For each such Event ID, it then sends
an Identify Producer message and checks:

\begin{enumerate}
\item That a Producer Identified or Producer Range Identified reply is
    received for each Identify Producer message sent.
\item That the reply carries the same Event ID as the request (or a range
    that encompasses it).
\end{enumerate}

Other messages received during this check are ignored.

\subsection{Track-Empty Sentinel (Informational)}

This section checks whether the DBC advertises Event IDs using the
track-empty sentinel value (0x3800 in the 14-bit address field).

Per the Working Note:
\textit{``A detector may, but does not have to, use the 0x3800
`track is empty' sentinel.''}

This check is informational only and never fails.  It reports:

\begin{enumerate}
\item Whether any detector Event IDs with address field 0x3800 are
    present among the Producer Identified replies.
\item If present, which direction/occupancy variants (arrives, leaves)
    are advertised.
\end{enumerate}

A detector that does not advertise track-empty events passes this check
with an informational note.

\end{document}
